\documentclass{article}
\usepackage{enumitem}

\title{Deep Reinforcement Learning Concepts Project}
\author{B9AI105 REINFORCEMENT LEARNING}
\date{18/3/2024}

\begin{document}

\maketitle

\section*{Objective}
Develop a comprehensive understanding of key concepts in deep reinforcement learning (DRL) and algorithms used in DRL, with a practical implementation in the Frozen Lake environment.

\section*{Task}
Choose one of the following topics related to DRL:
\begin{enumerate}[label=\arabic*.]
    \item Deep Q-Network (DQN)
    \item Replay Buffer
    \item Target Network
    \item Double DQN
    \item Prioritized Experience Replay
    \item Advantage Function
    \item Long Short-Term Memory (LSTM) in Deep Recurrent Q-Networks (DRQN)
    \item Value-Based Methods
    \item Policy-Based Methods
    \item Policy Gradient Method
    \item Reward-to-Go
    \item Policy Gradient with Baseline Function
    \item Baseline Function
\end{enumerate}

\subsection*{Implementation Task}
\begin{itemize}
    \item Implement the chosen concept or algorithm in the Frozen Lake environment using OpenAI Gym.
    \item Train an agent to navigate the Frozen Lake using the implemented algorithm.
    \item Compare the performance of your agent with different hyperparameters and settings.
\end{itemize}

\section*{Requirements}
\begin{enumerate}[label=\arabic*.]
    \item Write a detailed explanation of the chosen topic, covering the following aspects:
    \begin{itemize}
        \item Definition and purpose of the concept or algorithm.
        \item Why it is important or necessary in reinforcement learning.
        \item How it is used or implemented in DRL algorithms.
        \item Any advantages, disadvantages, or variations of the concept or algorithm.
    \end{itemize}
    \item Provide a code example or pseudocode illustrating the concept or algorithm in action in the Frozen Lake environment.
    \item Include references to academic papers, articles, or online resources that further explain the concept or algorithm.
\end{enumerate}

\section*{Deliverables}
\begin{itemize}
    \item A written report or presentation summarizing the chosen topic, including the explanation, code example (if applicable), and references.
    \item Code files for the implementation in the Frozen Lake environment.
\end{itemize}

\section*{Submission Guidelines}
\begin{itemize}
    \item Submit the report or presentation slides in PDF format.
    \item Include the code files in a separate folder or archive.
    \item Make sure to submit this assignment by \textbf{18th of April 2024}
\end{itemize}

\section*{Evaluation Criteria (Total: 100 points)}
\begin{enumerate}[label=\arabic*.]
    \item Explanation of Chosen Topic (30 points)
    \item Code Implementation (30 points)
    \item Performance Evaluation (20 points)
    \item Documentation (10 points)
    \item References and Further Reading (5 points)
    \item Overall Presentation (5 points)
\end{enumerate}

\end{document}



\section*{Learning Objectives}
This assessment aligns with the following Minimum Intended Module Learning Outcomes (MIMLOs) from the Reinforcement Learning module:
\begin{itemize}
  \item \textbf{MIMLO 7.2:} Propose Reinforcement Learning methods for complex problem scenarios.
  \item \textbf{MIMLO 7.3:} Apply appropriate RL techniques and algorithms to defined tasks.
  \item \textbf{MIMLO 7.4:} Implement RL frameworks in real-world or simulated problem environments.
\end{itemize}

\section*{Assessment Description}
Working in groups of 3–4, students are required to select, design, and implement a Reinforcement Learning solution to a defined problem. The problem domain is open and may be selected from, but not limited to:
\begin{itemize}
  \item Classic control tasks (e.g., CartPole, MountainCar)
  \item Grid-world or maze navigation (e.g., custom environments)
  \item Game environments (e.g., Blackjack, Atari, Unity ML-Agents)
  \item Robotics simulations (e.g., OpenAI Gym, MuJoCo)
  \item Resource allocation and scheduling problems
  \item Healthcare or finance RL applications
\end{itemize}
Students are free to choose any appropriate environment for their RL project. This may include pre-existing environments from OpenAI Gym or similar platforms, or they may design and implement a custom environment relevant to their area of interest or research.

\section*{Deliverables}
Each group must submit:
\begin{itemize}
  \item A technical report (3,000–4,000 words) including:
    \begin{itemize}
      \item Problem background and objective
      \item Algorithm selection and implementation details
      \item Evaluation metrics and analysis
      \item Reflections on ethical, technical, and practical challenges
    \end{itemize}
  \item Source code (via GitHub or archive)
  \item Meeting notes and individual contributions
\end{itemize}

\section*{Assessment Criteria}
\begin{tabular}{|p{8cm}|p{3cm}|}
\hline
\textbf{Criterion} & \textbf{Weighting} \\
\hline
Justification of RL technique & 20\% \\
\hline
Implementation quality & 20\% \\
\hline
Evaluation and performance analysis & 15\% \\
\hline
Reflection on challenges and ethics & 15\% \\
\hline
Technical writing and clarity & 10\% \\
\hline
Team collaboration evidence & 10\% \\
\hline
Relevance and problem significance & 10\% \\
\hline
\end{tabular}

\section*{Reassessment}
Reassessment will require an individual version of the above project. The problem should be scoped accordingly and the report must explicitly demonstrate how each MIMLO is addressed. Code and report must be submitted individually.
